\chapter{Introducción, objetivos y alcance académico}
\section{Introducción}
Diseñar una minicentral exige relacionar el recurso hídrico con la población atendida, el salto aprovechable y la capacidad real de las máquinas. En un sistema aislado, la frecuencia y la tensión dependen de sus propios generadores; el déficit de agua no puede compensarse importando del SEIN. El diseño debe incorporar regulación, respaldo o gestión de carga y explicar cómo cambia el servicio durante el año.

La instrucción del docente fue diseñar una minicentral hidroeléctrica, sin una rúbrica técnica detallada proporcionada. Este trabajo define un alcance académico integrado: investigación pública, modelación, selección argumentada y comprobaciones de ingeniería con sus límites. Las normas aplicables orientan las verificaciones; el estudio no reemplaza los permisos, los ensayos ni la responsabilidad profesional de una obra real.

El proyecto parte del borrador recibido y del trabajo previo de redes de distribución de Aquiles. R03 modifica el modo de operación que R02 había supuesto conectado a red. Retira el objetivo rígido de exportación de 150 kW, introduce población, caudales públicos simulados y una comparación de tamaños. El rango 150--200 kW se conserva como alternativa inicial, abierto a cambios justificados.

\section{Objetivo general}
Desarrollar un proyecto académico de diseño para replantear el aprovechamiento hidroeléctrico de Limbani y atender, mediante un sistema aislado, la demanda de Limbani y Huancasayani, identificando las restricciones del recurso, las condiciones de seguridad y la viabilidad de las alternativas.

\section{Objetivos específicos}
\begin{enumerate}
\item Identificar fuentes ANA, SENAMHI, INEI y antecedentes cartográficos; distinguir observación, simulación, estimación y decisión de diseño.
\item Analizar variación estacional, estiaje, reservas de agua, duración de caudales y disponibilidad de potencia/energía.
\item Construir una demanda por localidad y contrastarla con tamaños de planta, sin duplicar el total distrital ni usar huellas sintéticas como usuarios reales.
\item Dimensionar obras hidráulicas, grupo, tubería, auxiliares y sistema eléctrico; comprobar pérdidas y transitorios con hipótesis trazables.
\item Diseñar la regulación y las funciones de protección para operación aislada y describir previsiones para una eventual interconexión futura.
\item Elaborar metrados, costos, planos académicos y conclusiones coherentes con los cálculos y la matriz normativa.
\end{enumerate}

\section{Ficha de alcance}
\begin{longtable}{L{4.2cm}L{10cm}}
\toprule Elemento & Decisión y condición\\\midrule
Naturaleza & Proyecto académico estudiantil de diseño; replanteamiento de un aprovechamiento existente referido por el estudiante\\
Beneficiarios & Limbani y agrupación Huancasayani en el distrito de Limbani; Centro, Bellavista y Lince se distinguen por código censal\\
Operación inicial & Sistema aislado, 60 Hz; sin enlace ni importación del SEIN\\
Potencia & Comparación 150, 200, 400 y 600 kW; alternativa de cobertura conservadora 3 por 200 kW, pendiente de padrón y aforos\\
Obras & Derivación de pasada, desarenado, canal, cámara, tubería, casa de máquinas y restitución\\
Electricidad & 380 V de generación, transformación a 10 kV y distribución 380/220 V como propuesta académica coherente con el antecedente\\
Continuidad & Hidro más respaldo y prioridad de cargas; no se atribuye suministro permanente a la placa hidroeléctrica\\
Horizonte & Diseño de escenario 2026, sensibilidad de crecimiento y previsión de eventual SEIN entre 2046 y 2076\\
Entregables & Memoria y anexos, planos/esquemas académicos, tablas reproducibles y material de exposición\\
\bottomrule\end{longtable}

\section{Método y trazabilidad}
El directorio de datos guarda entradas con unidades, procedencia y estado. Los originales recibidos y los archivos descargados se conservan inalterados. Los resultados se generan a partir de esas entradas; las tablas no se ajustan manualmente para producir una potencia deseada. Las cantidades calculadas alimentan el presupuesto y los esquemas.

\begingroup\small
\begin{longtable}{L{3.2cm}L{6cm}L{4.8cm}}
\toprule Tipo de información & Ejemplo del trabajo & Alcance\\\midrule
Dato documental & Población INEI 2017, coordenadas de una ficha ANA & Verificado en su fuente; no se convierte en dato actual o de otra instalación\\
Simulación pública & Caudales PISCO ARNOVIC y GR2M & Producto modelado; no registro de aforos\\
Estimación del proyecto & Transposición por área, demanda proyectada, costos & Método y supuestos expresos; calibración pendiente\\
Decisión académica & Configuración aislada, diámetros, equipos y regulación & Propuesta para comparar y comprobar\\
Dato no disponible & Padrón activo, resistividad, geotecnia y placa existente & Se registra como pendiente; no se sustituye por una medición ficticia\\
\bottomrule\end{longtable}\endgroup
